\section{Method: Conjunctive Task-Circuit Saliency}
\label{sec:method}

We keep the \baseline{} mixed-precision quantization pipeline (GPTQ with
per-weight bit-width mask) and replace only the saliency extraction stage.

\paragraph{Setup.}
Model $f_\theta$ with weights $\theta$; uniform $b$-bit reference
$W_{b\text{bit}}$ (RTN or GPTQ on task calibration data); outlier budget $k$
(fraction of parameters kept in FP16, 0.35\% throughout). Contrastive pairs
$(x^+, x^-)$ where $x^+$ is task-faithful and $x^-$ is corrupted (schema
errors, wrong CoT, wrong answer choice, etc.).

\paragraph{\baseline{} saliency (review).}
Accumulate $|\nabla_W \mathcal{L}_{\text{CE}}(x^+)|$ over calibration, multiply
by $|W| \cdot |W - W_{b\text{bit}}|$, global top-$k$
selection~\citep{xiao2025tacq}. Gradients are taken w.r.t.\ polysemantic
weight tensors: each $W$ may contribute to many superposed features, and
$\mathcal{L}_{\text{CE}}$ mixes all roles active at each token.

\subsection{Circuit saliency in a feature basis}
\label{sec:method-circuit}

The circuit term inserts a feature-space bottleneck before weight saliency.
For each layer $\ell$, an encoder maps the MLP computation to a code
$f \in \mathbb{R}^{F_\ell}$; a decoder $D_\ell$ maps codes back to the MLP
output space. We rank \emph{task-discriminative} features via contrastive
activation: on corrupted input $x^-$ (no grad) record
$a^{\text{corr}}_{\ell,f} = \max_t f_{\ell,t,f}$; on clean input $x^+$ (with
grad) compute $a^{\text{clean}}_{\ell,f}$ analogously. Task features satisfy
\begin{equation}
\Delta a_\ell \;=\; \mathrm{ReLU}\!\left(a^{\text{clean}}_\ell - a^{\text{corr}}_\ell\right) \;>\; \tau_\ell,
\end{equation}
with $\tau_\ell$ the per-example 95th percentile of positive $\Delta a$. The
target MLP output contribution is
$y^{\text{target}}_\ell = D_\ell\!\left(f^{\text{clean}}_\ell \odot \mathbf{1}_{\Delta a > \tau}\right)$,
and the per-layer alignment loss is
$\mathcal{L}_\ell = -\sum_t \langle h^{\text{mlp,out}}_{\ell,t},\, y^{\text{target}}_{\ell,t} \rangle$.
One backward pass over $\mathcal{L} = \sum_\ell \mathcal{L}_\ell$ accumulates
$|\nabla_W \mathcal{L}|$ (\texttt{sample\_abs} semantics, matching
\baseline{}). Only weights on paths that implement the selected
task-discriminative features receive gradient; polysemantic CE would
additionally reward weights tied to non-task features active at the same
tokens.

\paragraph{Basis 1: transcoders (\method{}).}
$F_\ell \gg d_{\text{mlp}}$ sparse TopK transcoder features
(\texttt{facebook/crv-8b-instruct-transcoders}; $F{=}131{,}072$), whose entries
are approximately
monosemantic~\citep{bricken2023monosemanticity,dunefsky2024transcoders}. This
basis offers the cleanest feature separation but exists only for models with a
published dictionary.

\paragraph{Basis 2: native MLP neurons (\methodfree{}, dictionary-free).}
\label{sec:method-dictfree}
We take the model's own MLP hidden layer as the feature basis: the
post-activation intermediate
$f_{\ell} = \sigma(W_{\text{gate}} x) \odot (W_{\text{up}} x) \in \mathbb{R}^{d_{\text{mlp}}}$
serves as the code, and the model's own \texttt{down\_proj} matrix is the
decoder $D_\ell$. The contrastive selection rule and alignment loss are
\emph{identical} to the transcoder case; only the basis changes.

\emph{Intuition.} Individual MLP neurons are more polysemantic than dictionary
features, so each selected neuron is a noisier task indicator. But the
contrastive $\Delta a$ filter operates at the \emph{population} level: neurons
that consistently fire more on clean than corrupted task inputs are, in
aggregate, enriched for the task circuit even if no single neuron is
monosemantic. Because the conjunction with CE saliency (below) demotes weights
that only one signal favors, residual polysemantic noise in the circuit term
is filtered rather than propagated into the mask. Empirically this native
basis recovers most of the transcoder's gain on GSM8K and surpasses it on
MMLU-STEM (\S\ref{sec:ablation})---the dictionary is an optional booster, not a
requirement. \methodfree{} thus applies to \emph{any} transformer with no
external artifacts, removing the main practical barrier to task-circuit PTQ.

\subsection{Conjunctive combination and mask}
\label{sec:method-conjunction}

For each weight tensor $W$, define per-weight CE saliency $g_{\text{CE}}$ as in
\baseline{} and circuit saliency $g_{\text{align}}$ from $\mathcal{L}$ above
(same $|W|$ and $|W - W_{b\text{bit}}|$ scaling). Our main variant
(\texttt{mult}) combines both via geometric mean:
\begin{equation}
\label{eq:mult}
S(W) \;=\; |W| \cdot |W - W_{b\text{bit}}| \cdot \sqrt{g_{\text{CE}}(W)\, g_{\text{align}}(W)},
\end{equation}
selecting weights salient under \emph{both} task loss and task circuit. The
conjunction is the active ingredient: replacing $g_{\text{CE}}$ entirely with
$g_{\text{align}}$ (\texttt{align}) or using CE alone ties or trails on every
task we test, and union-style boosts
($g_{\text{CE}}(1 + \lambda g_{\text{align}})$) degenerate at 2-bit
(\S\ref{sec:ablation}). Global top-$k$ over all eligible linear weights yields
the binary mask $M$; GPTQ quantizes $M(W){=}0$ to $b$-bit and keeps
$M(W){=}1$ in FP16.

\paragraph{Why the conjunction works at low bit-width.}
At 2-bit, quantization noise is large enough that protecting the wrong weights
has a first-order cost: budget spent on weights that merely \emph{encode}
generic surface statistics (high CE saliency, low circuit saliency) or on
circuit edges whose quantization error happens to be benign (high circuit, low
CE) is budget taken from weights that are both functionally necessary and
fragile. The geometric mean implements an AND-gate over the two failure modes.
At 3-bit the noise floor drops and nearly any sensible task-conditioned mask
suffices---which is exactly what we observe (\S\ref{sec:results-3bit}).

\subsection{Anchored union: consensus core and dispute region}
\label{sec:method-anchored}

Empirically, \baseline{} and v2~mult agree on only ${\sim}46$--$50\%$ of the
FP16 budget ($J{\approx}0.30$--$0.33$; \S\ref{sec:analysis-overlap}). We
formalize a two-region policy that never renegotiates consensus outliers and
allocates only the remaining budget on the method-disputed mass. Let $\Omega$
be the eligible weight indices, $s_{\text{\baseline{}}}(w)$ and
$s_{\text{\method{}}}(w)$ the final per-weight scores, and
$k = \lfloor \rho |\Omega| \rfloor$, $\rho{=}0.0035$. Each method induces a
mask by global top-$k$; define
$M_{\text{core}} \triangleq M_{\text{\baseline{}}} \cap M_{\text{\method{}}}$
and
$M_{\text{dispute}} \triangleq M_{\text{\baseline{}}} \,\triangle\, M_{\text{\method{}}}$.
Stage~1 locks the core ($M \supseteq M_{\text{core}}$); Stage~2 fills the
residual budget $k_{\text{fill}} = k - |M_{\text{core}}|$ from
$M_{\text{dispute}}$ by ranking the blended normalized score
$s_\lambda(w) = \lambda\, \hat{s}_{\text{\method{}}}(w) + (1-\lambda)\, \hat{s}_{\text{\baseline{}}}(w)$.
The sweep over $\lambda$ tests where the optimal dispute allocation lies as
quantization severity changes (\S\ref{sec:results}).

\paragraph{Implementation notes.}
Transcoders (when used) stream from CPU; saliency ranks on CPU. Native-basis
extraction adds \emph{no} memory overhead beyond hooks on the MLP intermediate.
Model: Llama-3.1-8B-Instruct. Complexity matches \baseline{} asymptotics (one
forward--backward per calibration example), plus per-layer encode/decode for
the transcoder basis only. \todo{Report wall-clock vs.\ \baseline{} on 128
pairs; add Algorithm~1 pseudocode.}
